\section{Discussion and limitations}\label{sec:discussion}

%% ⛔ PROSA SUBSTITUIDA em 04/10/2026 pelo texto do Prof. Rudinei (e-mail ao Jader): secao
%% corrida, sem a subsecao Limitations, e sem parte das autocriticas, que ele tirou de
%% proposito («ou vao confundir o revisor ou ja sao esperadas»). NAO restaurar as autocriticas.
%% Aplicada com quatro ajustes, decididos pelo Jader em 04/10, que mudam o que o texto AFIRMA:
%%  1. «keypoint density» -> «descriptor count»: a fatoracao e em D_i, e keypoint density e
%%     outro controle diagnostico do Metodo (sec_method, controles diagnosticos).
%%  2. o ruido nas OUTRAS operacoes fica como consistencia, nao causa: o mecanismo so e medido
%%     no SIFT+BoVW, e a Conclusao diz «measured for that score and consistent with the others».
%%  3. a citacao do Otani e o ponteiro as Tabelas S24 e S25 voltaram.
%%  4. sem «across hardware runs» (causa que o artigo nao afirma) e «later», nao «prospective».
%%  5. «visual quality» saiu duas vezes (convencao de termos de 08/09).
%% Os numeros que ja eram macro continuam macro.

The central finding of this study is that visual-selection statistics derived from the video summarization literature, when reused as standalone technical-quality scores, fail to satisfy the technical-fidelity requirements proposed here. The algebraic factorization of the SIFT+BoVW raw-count score demonstrates that its behavior is overwhelmingly driven by descriptor count rather than visual fidelity; replacing the score with the raw descriptor count preserves every single choice across the \redDe{} operational clusters. For other operations, the concentration of reversals under white noise is consistent with noise inflating local high-frequency detail, which these statistics can misread as superior quality. At the application level, our evaluation separates scoring behavior from output preference, establishing selection sensitivity rather than guaranteeing product improvement. Consequently, practical reuses of these operations require explicit documentation of the intended property, exact statistics, declared polarity, sampling, and aggregation regimes. Although dedicated no-reference image quality assessment (NR-IQA) models offer better-supported alternatives, their robustness remains tightly bound to training exposure, distortion types, and domain characteristics. These findings reframe two readings in the current literature: first, extending the evaluation-side insight of Ref.~\citep{otani2019rethinking}, that protocol choices can dominate summarization scores, to the selection stage; and second, showing that isolated visual attributes frequently fail when extracted from the broader objective functions in which they were originally embedded.

These conclusions must be interpreted alongside some inherent limitations. First, while the quantitative measurements are fully deterministic given the benchmark inventory, the upstream process of compiling the archive, coding each source's functional role, and recording validation evidence relied on manual analysis by a single author. To ensure reproducibility and verification, all \valComPdf{} source works were read in full, and every attribution is anchored to recorded verbatim excerpts (Tables~S24 and~S25). Furthermore, our experimental scope is restricted to selected implementations across three image-quality databases and one broadcast-news pipeline; as a result, these conclusions do not automatically extend to different tasks, model families, or newer neural checkpoints.

The design of the benchmark and inference protocols also introduces specific constraints. The constructed evaluation benchmark comprises twenty video segments organized into \invBlocosN{} source-aware blocks, where ground-truth fidelity is established via controlled algorithmic transformations rather than subjective human preference. Statistical inference remains sensitive to small-sample tail conventions and spatial dependence assumptions, particularly given that eligibility thresholds and most inference parameters were not preregistered. Consequently, analytical variations such as best-component contrasts, oracle polarity alignment, and in-sample ridge regressions serve purely as diagnostic insights rather than prospective validations.

At the operational application level, the 42 redundancy clusters represent operationally redundant segment alternatives nested across 19 news topics with shared source material, rather than strictly semantically equivalent recordings. Because no human subjective quality ratings were collected for these specific segments, the reference benchmarks evaluate response to controlled signal degradation rather than human preference between distinct video captures. Technical quality here functions strictly as a tie-breaker within redundancy groups rather than a holistic verdict on the resulting summary. Similarly, external DOVER comparisons cannot isolate temporal modeling effects from underlying neural architecture or training data, and the official standardized fusion model was omitted.

Finally, the auxiliary diagnostic analyses presented are descriptive rather than confirmatory. In particular, conditioning the SIFT+BoVW score directly on keypoint counts risks removing legitimate quality signals alongside the dominant count channel. While the archived execution outputs allow full reconstruction of all derived statistics, the original software execution environment is not bit-identical. Lastly, partial-coverage models and later deep control architectures remain outside the primary screening comparisons evaluated in this audit.
